{\footnotesize
\begin{longtable}{@{}>{\raggedright\arraybackslash}p{0.21\textwidth} l c r r >{\raggedright\arraybackslash}p{0.54\textwidth}@{}}
\caption{All 49 enabled rules, generated from the rule metadata (\texttt{python3 -m tsw --list-rules}).
Type: hit, rate (Poisson), choice (binomial), stat (Gaussian), pred.\ (predicate).
Rules are named \texttt{family.rule}; the family prefix is given by the group headings. Dir: direction of a hit (for choice rules, of variant A). Cap: per-rule cap on $|w \cdot \mathrm{LLR}|$; w: weight.}
\label{tab-appendix}\\
\toprule
rule & type & dir & cap & w & description \\
\midrule
\endfirsthead
\toprule
rule & type & dir & cap & w & description \\
\midrule
\endhead
\addlinespace[3pt]\multicolumn{6}{@{}l}{\textit{starter} (2 rules, family cap 6)}\\*
\texttt{modal\_particles} & rate & O & 3 & 1.00 & modal particles ju/nog/väl/visst \\
\texttt{sent\_len} & stat & T & 2 & 0.50 & words per sentence \\
\addlinespace[3pt]\multicolumn{6}{@{}l}{\textit{connectives \& punctuation} (13 rules, family cap 4)}\\*
\texttt{native\_connectives} & rate & O & 2 & 0.50 & native connective adverbs (dessutom, därför, dock, därmed, inte minst \ldots{}) \\
\texttt{additive\_ocksa} & rate & O & 1.5 & 0.40 & additive `också' \\
\texttt{calque\_connectives} & rate & T & 2 & 0.80 & Swedish equivalents of English discourse markers (i synnerhet, i själva verket \ldots{}) \\
\texttt{att\_drop} & rate & O & 1.5 & 1.00 & complementizer `att' dropped after tror/tycker/anser/hoppas/säger \ldots{} \\
\texttt{som\_drop} & choice & T & 1.5 & 0.70 & object relative without `som' (`det arbete ni har gjort') \\
\texttt{initial\_och} & rate & T & 2 & 1.00 & sentence-initial `Och' \\
\texttt{initial\_men} & rate & T & 1.5 & 0.70 & sentence-initial `Men' \\
\texttt{initial\_adverbial\_comma} & choice & T & 2 & 1.00 & comma after a sentence-initial adverbial (`Dessutom, \ldots{}') \\
\texttt{clause\_comma\_omission} & choice & T & 1.5 & 0.70 & no comma before clausal men/eftersom/medan/vilket \\
\texttt{written\_out\_abbrev} & choice & T & 1.5 & 0.50 & `till exempel/så kallad/till och med' written out vs `t.ex./s.k./t.o.m.' \\
\texttt{parenthetical\_dash} & rate & T & 1.5 & 0.50 & dash between words (parenthetical insertion) \\
\texttt{explanatory\_parens} & rate & T & 1 & 1.00 & explanatory parentheses (no numbers, not document codes) \\
\texttt{exclamation} & pred. & O & 1.5 & 0.50 & exclamation marks outside vocatives \\
\addlinespace[3pt]\multicolumn{6}{@{}l}{\textit{register} (9 rules, family cap 6)}\\*
\texttt{kommer\_att\_vs\_ska} & choice & T & 2.5 & 1.00 & future/intention auxiliary `kommer att' vs `ska/skall' \\
\texttt{skulle\_vilja\_vs\_vill} & choice & T & 2 & 1.00 & `jag/vi skulle vilja' vs `jag/vi vill' \\
\texttt{tycker\_vs\_anser} & choice & O & 2 & 1.00 & opinion verb `tycka' vs `anse' \\
\texttt{valdigt\_vs\_mycket} & choice & O & 1.5 & 1.00 & degree adverb `väldigt' vs `mycket' + adjective \\
\texttt{sjalvklart} & rate & O & 1.5 & 1.00 & stance adverb `självklart' \\
\texttt{dem\_har\_rate} & rate & T & 1 & 0.50 & demonstrative `den/det/de här' \\
\texttt{bland\_annat} & rate & O & 1.5 & 0.60 & `bland annat / bl.a.' \\
\texttt{samt} & rate & O & 1.5 & 1.00 & coordinator `samt' \\
\texttt{bra} & rate & O & 1.5 & 1.00 & evaluative `bra' \\
\addlinespace[3pt]\multicolumn{6}{@{}l}{\textit{particles \& pronouns} (15 rules, family cap 4)}\\*
\texttt{namligen} & rate & O & 2.5 & 1.00 & clause-internal particle `nämligen' (explanatory `you see') \\
\texttt{anda} & rate & O & 2 & 1.00 & concessive particle `ändå' \\
\texttt{alltsa} & rate & O & 1.5 & 1.00 & inferential particle `alltså' \\
\texttt{stance} & rate & O & 2 & 1.00 & stance adverbs `egentligen/visserligen/i och för sig/i alla fall' \\
\texttt{da\_anaphoric} & rate & O & 2 & 1.00 & anaphoric/inferential `då' (`then, in that case') \\
\texttt{of\_course} & choice & T & 1.5 & 1.00 & `of course': `naturligtvis' vs `givetvis/förstås' \\
\texttt{dvs} & rate & O & 1.5 & 1.00 & reformulation marker `dvs.' / `det vill säga' \\
\texttt{particle\_verbs} & rate & O & 1.5 & 1.00 & Swedish-specific particle verbs (slå fast, ställa sig bakom, lyfta fram, \ldots{}) \\
\texttt{ta\_itu\_med} & rate & T & 2 & 1.00 & `ta itu med' (EN tackle/address/deal with) \\
\texttt{generic\_man} & rate & O & 1.5 & 0.25 & generic pronoun `man' \\
\texttt{copular\_det} & rate & O & 1.5 & 0.50 & copular `det är/var' and `är/var det' (expletive, cleft, anaphoric det) \\
\texttt{poss\_1st} & rate & T & 1.5 & 0.60 & 1st-person possessive determiners `min/mitt/mina/vår/vårt/våra' \\
\texttt{poss\_3rd} & rate & T & 2 & 0.80 & 3rd-person non-reflexive possessives `hans/hennes/deras/dess' \\
\texttt{han\_hon} & rate & T & 1 & 0.50 & 3rd-person personal pronouns `han/hon' \\
\texttt{ni\_vs\_man} & choice & T & 1.5 & 1.00 & `ni' vs generic `man' \\
\addlinespace[3pt]\multicolumn{6}{@{}l}{\textit{calques} (4 rules, family cap 3)}\\*
\texttt{lexicon} & hit & T & 3 & 1.00 & English calques and anglicisms \\
\texttt{native\_counterparts} & hit & O & 2 & 1.00 & native alternatives of calques \\
\texttt{forsoka\_att} & choice & T & 1.5 & 1.00 & `försöka att + V' vs `försöka + V' \\
\texttt{pied\_piped\_relative} & hit & T & 1.5 & 1.00 & pied-piped relative `P + vilken/vilket/vilka' \\
\addlinespace[3pt]\multicolumn{6}{@{}l}{\textit{word order} (3 rules, family cap 4)}\\*
\texttt{fronted\_theme} & pred. & O & 3 & 0.80 & topicalised non-subject in the prefield (V2 inversion): which kind \\
\texttt{subject\_theme} & pred. & T & 2 & 0.80 & subject in the prefield: pronoun / NP / det / demonstrative / cleft \\
\texttt{v2\_violation} & pred. & T & 3 & 1.00 & English word order leaking into Swedish (V3 / no inversion / embedded negation) \\
\addlinespace[3pt]\multicolumn{6}{@{}l}{\textit{statistics} (3 rules, family cap 2)}\\*
\texttt{long\_words} & stat & T & 1 & 0.50 & share of long words (> 6 letters, the LIX criterion) \\
\texttt{av\_nominal} & stat & T & 1 & 1.00 & `av' after a definite nominalisation (`genomförandet av') \\
\texttt{av\_other} & stat & T & 1.5 & 1.00 & `av' elsewhere (of-periphrasis, by-agent) \\
\bottomrule
\end{longtable}}
